\section{Related Work}
\noindent\textbf{Large activations in transformers:} 
% {\textcolor{blue}{Can be removed for space IMO-Understanding the behavior of transformers has been an important research topic in the last few years, owing to their success in various NLP and vision tasks. A common approach employed to achieve this is to analyze the activations and the features extracted by the model.}}
Understanding the behavior of transformers has been an important research topic in the last few years, owing to their success in various NLP and vision tasks. A common approach employed to achieve this is to analyze the activations and the features extracted by the model. Recent studies have uncovered the presence of outlier weights and large-magnitude activations in transformers, affecting both language \cite{heimersheim2023residual,kovaleva2021bert,zhao2023unveiling,timkey2021all} and vision models \cite{darcet2024vision,mambaregisters}. These phenomena manifest as concentrated attention weights in LLMs \cite{kovaleva2021bert,zhao2023unveiling} and in vision models~\cite{darcet2024vision}

\noindent\textbf{Implications and solutions for massive activations:} These large activations are found to pose challenges for model quantization \cite{dettmers2022gptint,he2024understanding} and image segmentation \cite{darcet2024vision}. Proposed remedies to address these challenges include appending ``register'' tokens to the input patch sequence \cite{darcet2024vision} or explicit attention biases \cite{sun2024massive} during training. While these approaches typically discard the additional tokens during inference, our work identifies that these tokens contain valuable auxiliary information to that captured by the \texttt{[CLS]} token. Therefore, we propose leveraging this information to obtain richer feature representations, enabling superior performance across a diverse range of tasks.


% Our approach also differs from approaches presented in \cite{mambaregisters} that use the additional tokens solely for classifier improvement during training.


% \noindent\textbf{Presence of activations and features with large magnitudes in LLMs and ViTs}
% Understanding the behavior of transformers has been an important research topic in the last few years, owing to their success in various NLP and vision tasks. A common approach employed to achieve this is to analyze the activations and the features extracted by the model. Many works have uncovered intriguing properties such as the presence of outlier weights in the LayerNorm of GPT-2 and LLaMA2-13B \cite{heimersheim2023residual},\cite{kovaleva2021bert}, \cite{zhao2023unveiling} and the existence of feature dimensions containing activations with large magnitudes and found that these few ``rogue" dimensions dominate measures such as cosine similarity for evaluating representation similarity \cite{timkey2021all, heimersheim2023residual}. 
% Similarly, many works have also found the presence of artifacts in the attention maps of models where the attention weights are concentrated on few tokens such as the `SEP' token in the BERT models~\cite{kovaleva2021bert,zhao2023unveiling}. Darcet et al.~\cite
% {darcet2024vision} have shown that the attention maps of sufficiently large vision transformers are highly concentrated on few tokens and that these tokens correspond to the low-informative background tokens in the image. Recent efforts by Sun et al.~\cite{sun2024massive} sheds light on why these artificats in attention maps emerge and they show that these are due to the presence of massive activations that act as constant implicit biases in the self-attention mechanisms.  Recently, these activations were also found to be present in architectures such as Vision Mamba~\cite{mambaregisters}

% \noindent\textbf{Implications and remedies for these large activations} Due to the widespread presence of these large activations in models, many works have focused on their implications on downstream tasks. For instance, it has been shown that these activations and thus the outlier features pose significant challenges for effective quantization of the transformers~\cite{dettmers2022gptint,he2024understanding}, or in the case of ViTs, these activations pose challenges for processing the attention maps to segment objects in an image~\cite{darcet2024vision}. Due to these undesirable properties, Darcet et al.~\cite{darcet2024vision} have proposed to train ViTs with additional tokens appended to the input sequence and these `registers' are then found to inherit the propreties of the outlier tokens and have thus resulted in enhanced performance on segmentation tasks. Similarly, Sun et al.~\cite{sun2024massive} have proposed to eliminate the massive activations augmenting LLMs with explicit attention biases. 

% Importantly, while in both these cases, the additional tokens are discarded in the inference stage and we on the other hand, propose to leverage the information captured by these additional tokens for enhancing the performance of models on various tasks. Our proposed approach is also in contrast to the recent effort by Wang et al.~\cite{mambaregisters} that proposes to utilitze the additional tokens during training of the Vision Mamba models to obtain a better classifier. 


% Simulatenously, many works have also focused on the the implications 